% © 2025 Ing. David Jaroš — CC BY-NC-ND 4.0
%
% This work is licensed under a Creative Commons Attribution-NonCommercial-NoDerivatives 
% 4.0 International License (CC BY-NC-ND 4.0).
%
% License History: Earlier drafts (up to v0.3) were released under CC BY 4.0. 
% From v0.4 onward, all material is released under CC BY-NC-ND 4.0 to protect 
% the integrity of the theoretical work during ongoing academic development.
%
% See LICENSE.md for full license text.

\ifdefined\INCLUDEMODE
\else
  \documentclass[12pt]{article}
  \usepackage[a4paper, margin=2.5cm]{geometry}
  \usepackage{amsmath, amssymb, amsthm}
  \usepackage{hyperref}
  \usepackage{xcolor}
  \usepackage{mdframed}

  \newtheorem{theorem}{Theorem}[section]
  \newtheorem{lemma}[theorem]{Lemma}
  \newtheorem{proposition}[theorem]{Proposition}
  \theoremstyle{definition}
  \newtheorem{definition}[theorem]{Definition}
  \theoremstyle{remark}
  \newtheorem{remark}[theorem]{Remark}

  \newmdenv[backgroundcolor=yellow!20, linecolor=orange, linewidth=1.5pt]{gapbox}
  \newmdenv[backgroundcolor=green!15, linecolor=green!60!black, linewidth=1.5pt]{resultbox}

  \title{\textbf{Step 2: Electron Mass from Yukawa Coupling at $\varphi = \mathrm{const}$}}
  \author{David Jaroš}
  \date{2026-03-06}

  \begin{document}
  \maketitle

  \begin{abstract}
  We derive the electron mass from the Yukawa-type coupling between the biquaternionic
  scalar background $\varphi = v = \mathrm{const}$ and the spinor field $\psi$.
  The coupling $\mathcal{L}_Y = y\,\varphi\,\bar\psi\psi$ gives $m_e = y\cdot v$ in exact
  analogy with the Standard Model Higgs mechanism.  We discuss which aspects are derived
  from UBT structure and which remain as free parameters at the current stage.
  \textbf{Verdict: SKETCH — structural derivation complete; $y$ and $v$ not yet derived
  from first principles.}
  \end{abstract}
\fi

\section{Step 2: Electron Mass from Yukawa Coupling at $\varphi = \mathrm{const}$}
\label{sec:electron_mass}

\textbf{Task reference:} UBT\_v29\_task7, step2\_electron\_mass.
\textbf{Date:} 2026-03-06.
\textbf{Verdict: SKETCH.}

\subsection{Dirac Structure in UBT}

From Step~3 of the FPE verification chain
(\texttt{consolidation\_project/FPE\_verification/step3\_dirac\_emergence.tex}),
the biquaternion algebra
$\mathbb{B} = \mathbb{C}\otimes\mathbb{H}\cong\mathrm{Mat}(2,\mathbb{C})$
provides a natural representation of the Dirac gamma matrices:
\begin{equation}
  \gamma^\mu \;\in\; \mathbb{B}\otimes\mathbb{B},
  \qquad \{\gamma^\mu,\gamma^\nu\} = 2\eta^{\mu\nu}.
\end{equation}
The spinor field $\psi(x)$ emerges as the restriction of $\Theta$ to the spinorial
subspace of $\mathbb{B}$.

\subsection{Yukawa Coupling from the Action}

The UBT action $S[\Theta]$ contains an interaction term between the charged spinorial
component and the neutral scalar zero-mode $\theta_0$ of $\Theta$ (see Appendix AA,
eq.~(A.7)):
\begin{equation}
  \mathcal{L}_Y = y\,\theta_0\,\bar\psi\psi,
\end{equation}
where $y$ is the dimensionless Yukawa coupling constant.

\begin{gapbox}
\textbf{Gap Y1 (open):} The Yukawa coupling $y$ is not yet derived from first principles
within UBT.  It depends on the projection of the biquaternionic interaction vertex onto
the spinorial subspace.  A complete derivation would require evaluating the matrix element
$\langle\mathrm{spinor}|\mathcal{L}_\mathrm{int}|\mathrm{scalar}\otimes\mathrm{spinor}\rangle$
directly from $S[\Theta]$.
\end{gapbox}

\subsection{Structural Form of the Yukawa Coupling (Sketch)}
\label{subsec:yukawa_sketch}

While a full derivation of $y$ remains open (Gap Y1), its structural form can be
motivated from the biquaternionic geometry.

The UBT action $S[\Theta]$ contains, in the charged sector, a term of the form:
\begin{equation}
  \mathcal{L}_\mathrm{int} = g_\psi\,\Theta^\dagger \cdot \varphi \cdot \Theta,
\end{equation}
where $g_\psi$ is the gauge coupling on the $\psi$-circle and $\varphi$ is the scalar
zero-mode field.  The Yukawa coupling $y$ is then the projection of this vertex onto
the spinorial subspace:
\begin{equation}
  y = \mathrm{Sc}\!\left[e_0^\dagger \cdot (\partial_\varphi S[\Theta]) \cdot e_0\right],
  \label{eq:y_structural}
\end{equation}
where $e_0$ is the vacuum mode (the zero-mode spinor).

In dimensional terms, $y \sim g_\psi \cdot R_\psi / R_\varphi$, where $R_\psi$ is the
$\psi$-circle radius (determining the compactification scale) and $R_\varphi$ is the
scalar field radius in field space.

\begin{gapbox}
\textbf{Gap Y1 (SKETCH):} The Yukawa coupling has the structural form
$y \sim g \cdot R_\psi / R_\varphi$ where $g$ is the gauge coupling.
A full derivation requires computing $\partial_\varphi S[\Theta]$ explicitly ---
pending Gap M1 (soliton mass computation).

Status: \textbf{Y1 — SKETCH}.
\end{gapbox}

\subsection{Mass Generation}

For $\theta_0 = \langle\theta_0\rangle = v = \mathrm{const}$ (the constant scalar background
from the problem statement), the Yukawa term becomes a mass term:
\begin{equation}
  \mathcal{L}_Y\big|_{\varphi=v} = (y\cdot v)\,\bar\psi\psi \;\equiv\; m_e\,\bar\psi\psi,
\end{equation}
with
\begin{equation}
  \boxed{m_e = y\cdot v.}
\end{equation}
This is structurally identical to the Standard Model Higgs mechanism with $v$ playing the
role of the Higgs VEV.

\subsection{What Is $v$ in UBT?}

The scalar VEV $v$ is determined by the vacuum structure of $\Theta$ on the $\psi$-circle.
The $\psi$-zero-mode is the constant mode $e^{i\cdot0\cdot\psi/R_\psi} = 1$; its
coefficient $v$ is set by minimising the effective potential on the $\psi$-circle.

At one-loop level (see Appendix CT for the two-loop baseline), the effective potential
for $\theta_0$ reads schematically:
\begin{equation}
  V_\mathrm{eff}(\theta_0)
  = -\mu^2|\theta_0|^2 + \lambda|\theta_0|^4 + \cdots,
\end{equation}
with $\mu^2, \lambda$ determined by loop corrections from the charged sector.  The
minimum at $|\theta_0| = v = \mu/\sqrt{2\lambda}$ sets the scale of $v$.

\begin{gapbox}
\textbf{Gap Y2 (partially derived):} The parameters $\mu$ and $\lambda$ in $V_\mathrm{eff}(\theta_0)$
are not yet computed from $S[\Theta]$ without external input.  In particular:
\begin{itemize}
  \item $\mu^2$ requires a mass-squared parameter in the $\psi$-sector action;
  \item $\lambda$ arises from the $|\theta_0|^4$ self-interaction vertex.
\end{itemize}
Both are in principle determined by the topology of the $\psi$-circle (radius $R_\psi$)
and the biquaternionic geometry, but the explicit computation is an open task.
\end{gapbox}

\subsection{VEV from the Effective Potential (Partial Derivation)}
\label{subsec:vev_partial}

From the one-loop effective potential $V_\mathrm{eff}(\theta_0) = -\mu^2|\theta_0|^2 + \lambda|\theta_0|^4$,
the minimum is at:
\begin{equation}
  v^2 = \frac{\mu^2}{2\lambda}.
  \label{eq:vev_minimum}
\end{equation}
The renormalisation group analysis from Step~3 (see \texttt{step3\_beta\_function.tex})
shows that the $\beta$-function for the self-coupling $\lambda$ has a UV fixed point,
indicating that $V_\mathrm{eff}$ is renormalisationally stable.  This means:
\begin{itemize}
  \item $\lambda > 0$: the potential is bounded below (no runaway vacuum),
  \item $\mu^2 > 0$: spontaneous symmetry breaking occurs,
  \item $v = \mu/\sqrt{2\lambda}$ is a stable minimum.
\end{itemize}

In terms of the $\psi$-circle geometry:
\begin{equation}
  \mu^2 \sim \frac{1}{R_\psi^2}, \qquad \lambda \sim g_\psi^2,
\end{equation}
so that $v \sim 1/(g_\psi R_\psi)$.  The value of $v$ is thus set by the compactification
radius $R_\psi$ and the gauge coupling $g_\psi$ of the $\psi$-circle.

\begin{resultbox}
\textbf{Gap Y2 — PARTIALLY DERIVED}:
$V_\mathrm{eff}(\theta_0)$ has a minimum at $v^2 = \mu^2/(2\lambda)$ (equation~\eqref{eq:vev_minimum}).
Renormalisational stability (from Step~3 $\beta$-function analysis) confirms this
minimum is stable.  The explicit computation of $\mu^2$ and $\lambda$ from $S[\Theta]$
remains open; their scaling $\mu^2 \sim R_\psi^{-2}$, $\lambda \sim g_\psi^2$ is
dimensionally motivated.

Status: \textbf{Y2 — PARTIALLY DERIVED}.
\end{resultbox}

\subsection{Lepton Mass Ratios}

If the three lepton generations are identified with the $\psi$-modes $\Theta_0, \Theta_1,
\Theta_2$ (see \texttt{research\_tracks/three\_generations/}), then the mass ratios
$m_e : m_\mu : m_\tau$ follow from the Yukawa couplings of each mode:
\begin{equation}
  m_\ell = y_\ell \cdot v,
  \qquad \ell = e, \mu, \tau.
\end{equation}
Whether $y_\ell$ can be derived from the eigenvalue spectrum of the Hecke operators
(the Hecke conjecture, see \texttt{research\_tracks/three\_generations/step6\_hecke\_matches.tex})
is an open question classified as \textbf{Supported by numerical evidence}.

\subsection{Summary}

\begin{center}
\begin{tabular}{ll}
\hline
Claim & Electron mass $m_e = y\cdot v$ from Yukawa at $\varphi = \mathrm{const}$ \\
Status & \textbf{SKETCH} \\
Structural mechanism & Identical to SM Higgs; natural in $\mathbb{C}\otimes\mathbb{H}$ \\
Gap Y1 & Yukawa coupling $y \sim g\cdot R_\psi/R_\varphi$ — structural form (SKETCH) \\
Gap Y2 & VEV $v = \mu/\sqrt{2\lambda}$ from stable minimum of $V_\mathrm{eff}$ (PARTIALLY DERIVED) \\
Layer & [L1] — one-loop effective potential \\
\hline
\end{tabular}
\end{center}

\begin{resultbox}
\textbf{Conclusion (Step 2, SKETCH):}
The electron mass $m_e = y\cdot v$ arises from the Yukawa coupling of the spinor field
to the constant scalar background $\varphi = v$.  The mechanism is structurally identical
to the SM Higgs mechanism and is naturally realised within $\mathbb{C}\otimes\mathbb{H}$.
Two gaps remain: (Y1) derivation of $y$ from $S[\Theta]$, and (Y2) derivation of $v$
from the $\psi$-circle effective potential.  Neither gap affects the validity of Steps 1,
3, or 4 at the leading-order level.
\end{resultbox}

\ifdefined\INCLUDEMODE\else
  \end{document}
\fi
